% =====================================================================
% Section: Abstract  (outline.md "Abstract", 150-250 words)
% Draft v4 (2026-07-06). Wording mirrors the honest framing used in the body:
%   DPRNet matches CATANet-level accuracy with small gains on most benchmark cells
%   (1_introduction.tex, 3_experiments_comparison.tex, 3_experiments_efficiency.tex).
%   Numbers VERBATIM from paper/tables/* and paper/data/*: x4 Urban100 26.89 dB,
%   52.14G Multi-Adds (256x256-input protocol; below SwinIR-light 60.3G / SRFormer-light
%   56.5G, above CATANet 46.8G). Ablation is a 250k, 3-seed study: C1 dynamic
%   prototypes > EMA center, C4 balance loss -> more balanced + lower-variance usage.
%   No overstated compute/isolated-gain claims; consistent with method-experiment-traceability.md.
% =====================================================================
\begin{abstract}
Lightweight single image super-resolution (SR) must balance reconstruction quality with
model cost on resource-constrained platforms. Content-aware routing improves this balance
by grouping semantically related tokens, but the Token Aggregation Block (TAB) of CATANet
uses cross-batch exponential-moving-average centers, hard assignment labels, and a coupled
prototype update/confirmation step. We propose \emph{Dynamic Prototype Routing} (DPR), a
drop-in TAB replacement that keeps the CATANet backbone and tensor interface unchanged.
DPR forms per-image prototypes from the current tokens, refines them with a learnable
query bank, matches tokens with a bounded learnable temperature, and carries assignment
confidence into token ordering, global gating, and a soft fallback; a balancing loss
discourages prototype-slot collapse. Trained on DIV2K and evaluated on five standard
benchmarks at $\times2/\times3/\times4$, DPRNet stays in the lightweight regime and
matches CATANet-level accuracy, with small gains on most benchmark cells. At $\times4$,
for example, it obtains $26.89$\,dB on Urban100 with $52.14$\,G Multi-Adds. It uses fewer
Multi-Adds than the lightweight Transformer baselines SwinIR-light and SRFormer-light
under the same input convention. Multi-seed ablations confirm that the
dynamic prototypes outperform a history-center variant with an exponential moving average
and that the balancing loss makes prototype usage both more balanced and more stable across
seeds, while the routing diagnostics remain interpretable with soft assignment confidence.
\end{abstract}
